\section*{Erd\H{o}s Problem 1119: pointwise-small families of entire functions}
\subsection*{Statement and set-theoretic status}
Let $\aleph_0<\mathfrak m<\mathfrak c$. If a set $\mathcal F$ of distinct entire functions assumes at most $\mathfrak m$ values at every point, must $|\mathcal F|\le\mathfrak m$? This is a question about distinct functions, not repeated labels for a single function.
We prove the affirmative answer when $\mathfrak m^+<\mathfrak c$, identify exactly why the elementary argument stops, and reconstruct the countable-valued counterexample under CH as an analytic model of the obstruction. The boundary case $\mathfrak m^+=\mathfrak c$ is independent by the cited forcing results. Their forcing constructions are not reconstructed, so that proof obligation remains unresolved in this attempt.

\subsection*{The strict cardinal gap}
Distinct entire functions coincide on a countable set: their difference has isolated zeros, with only finitely many in each compact disk. Suppose $|\mathcal F|>\mathfrak m$ and choose $\mathcal G\subseteq\mathcal F$ of size $\mathfrak m^+$. The union
\[
 Z=\bigcup_{\substack{f,g\in\mathcal G\\f\ne g}}\{z:f(z)=g(z)\}
\]
has cardinality at most $(\mathfrak m^+)^2\aleph_0=\mathfrak m^+$. If $\mathfrak m^+<\mathfrak c$, choose $z\notin Z$. All $\mathfrak m^+$ values there are distinct, a contradiction. Constant functions with constants in a set of size $\mathfrak m$ show sharpness of the resulting bound.
Also there are only $\mathfrak c$ entire functions: each is determined by its countable Taylor-coefficient sequence, and $\mathfrak c^{\aleph_0}=\mathfrak c$. Thus at the problematic successor boundary a counterexample necessarily has size $\mathfrak c$.

\subsection*{A countable dense-interpolation lemma}
Let $z_1,z_2,\ldots$ be distinct complex numbers and let each $D_j\subset\mathbb C$ be countable and dense. There are continuum many entire functions satisfying $f(z_j)\in D_j$ for every $j$.
Build a binary tree of polynomials. At a node with polynomial $p$ already satisfying the first $n-1$ constraints, let
\[
 Q_n(z)=\prod_{j<n}(z-z_j).
\]
For either of two distinct values $d\in D_n$ sufficiently close to $p(z_n)$, set
\[
 p_{\mathrm{new}}(z)=p(z)+\frac{d-p(z_n)}{Q_n(z_n)}Q_n(z).
\]
Density supplies two such values so close that the added polynomial has modulus at most $2^{-n}$ on $|z|\le n$. The earlier constraints are preserved exactly, while the new one is satisfied. Start with $p=0$ and the empty product at the first step.
Along every infinite binary branch these corrections converge uniformly on compact sets, defining an entire function. Once a value at $z_n$ is fixed, all later corrections vanish there. Distinct branches have different values at their first branching point and give distinct functions. There are $2^{\aleph_0}$ branches, proving the lemma. Finite lists of constraints can be extended arbitrarily to countable lists.

\subsection*{What CH makes possible}
Assume CH and enumerate the complex numbers as $(z_\alpha)_{\alpha<\omega_1}$. Fix one countable dense set $D\subset\mathbb C$. Inductively choose an entire $f_\alpha$ distinct from its predecessors and with
\[
 f_\alpha(z_\beta)\in D\qquad(\beta<\alpha).
\]
This is possible: $\alpha$ is countable, the interpolation lemma supplies continuum many eligible functions, and only countably many predecessors are excluded. The family has cardinality $\aleph_1$. At $z_\beta$, all functions with $\alpha>\beta$ take values in $D$, while the functions with $\alpha\le\beta$ add at most countably many values. Hence the family is pointwise countable.
This is the classical countable-valued Wetzel counterexample under CH. It is \emph{not} a counterexample with the problem's uncountable $\mathfrak m$: under CH there is no cardinal strictly between $\aleph_0$ and $\mathfrak c$. Its purpose is to expose the dense-interpolation mechanism and the countability needed at each stage.

\subsection*{The actual successor boundary}
Schilhan--Weinert, \emph{Wetzel families and the continuum}, \url{https://arxiv.org/abs/2310.19473v3}, state and prove consistency of a Wetzel family when $\mathfrak c=\aleph_2$. Such a family has size $\aleph_2$ and pointwise at most $\aleph_1$ values, giving a negative answer with $\mathfrak m=\aleph_1$.
Their introduction also states the complementary Kumar--Shelah result: in the side-by-side Cohen model there is no Wetzel family. With continuum $\aleph_2$, a counterexample to the present assertion would itself be a Wetzel family by the cardinal bound above, so that model gives the positive answer.
These are explicit external relative-consistency inputs; no absolute contradiction in ZFC follows. At $\mathfrak m^+=\mathfrak c$, the coincidence union $Z$ can cover the plane, and the preceding elementary proof gives no point outside it. Repeating the CH interpolation over $\omega_2$ would require meeting uncountably many constraints, for which our lemma supplies no justification.

\subsection*{Assessment}
The strict-gap case, the interpolation lemma and the classical CH construction are proved. The independence classification is correctly explained from the primary literature, but its forcing arguments are not supplied. This improves on the older local attempt's suggestion to settle the boundary absolutely: the correct missing task is to reconstruct the consistency proofs.
Source statement: \url{https://www.erdosproblems.com/1119}. No novelty is claimed.
\subsection*{COMPLETION ESTIMATE}
\textbf{COMPLETION: 65\%}
